\documentclass[11pt]{article}
\usepackage[margin=1in]{geometry}
\usepackage{amsmath, amssymb, amsthm}
\usepackage{hyperref}

\newtheorem{theorem}{Theorem}
\newtheorem{proposition}[theorem]{Proposition}
\newtheorem{lemma}[theorem]{Lemma}
\newtheorem{corollary}[theorem]{Corollary}
\theoremstyle{remark}
\newtheorem{remark}[theorem]{Remark}

\DeclareMathOperator{\SYT}{SYT}
\newcommand{\Sym}{\mathrm{Sym}}
\newcommand{\ftilde}{\widetilde{f}}
\newcommand{\qe}{q_e}
\newcommand{\zetae}{\zeta_e}
\newcommand{\ZZ}{\mathbb{Z}}
\newcommand{\QQ}{\mathbb{Q}}
\newcommand{\CC}{\mathbb{C}}
\newcommand{\NN}{\mathbb{N}}
\newcommand{\one}{\mathbf{1}}

\title{Self-Similarity of the Fake-Degree Schur Sum at a Root of Unity}
\author{Clio}
\date{August 7, 2026}

\begin{document}
\maketitle

\begin{abstract}
For $n \ge 1$, $e \ge 2$, and $\zetae$ a primitive $e$-th root of unity, define
\[
  \qe^{(n)} \;:=\; \sum_{\lambda \vdash n} \ftilde_\lambda(\zetae)\, s_\lambda \;\in\; \Lambda_n \otimes \ZZ[\zetae],
\]
where $\ftilde_\lambda(t)$ is the fake-degree polynomial. We prove the self-similarity identity
\[
  \boxed{\;\qe^{(n)} \;=\; p_e^{\lfloor n/e\rfloor}\, \cdot\, \qe^{(n \bmod e)}\;}
\]
in the algebra of symmetric functions over $\ZZ[\zetae]$. In particular, the residue factor $F_e := \qe^{(n)}/p_e^{\lfloor n/e\rfloor}$ (a well-defined symmetric function of degree $n \bmod e$, by our earlier support theorem) depends on $n$ only through $r_0 := n \bmod e$, and its Schur expansion is given by the ``small'' partition-of-$r_0$ analogue:
\[
  F_e \;=\; \sum_{\lambda \vdash r_0} \ftilde_\lambda(\zetae)\, s_\lambda.
\]
The proof is a one-line consequence of the closed-form coefficient formula derived in \cite{Cli26-support}, using the classical identity $\prod_{r=1}^{e-1}(1 - \zetae^r) = e$.
\end{abstract}

\section{Statement and Context}

Fix $e \ge 2$ and let $\zetae := e^{2\pi i/e}$. For $n \ge 0$ set
\[
  \qe^{(n)} \;:=\; \sum_{\lambda \vdash n} \ftilde_\lambda(\zetae)\, s_\lambda,
\]
where the sum runs over partitions of $n$ and $\ftilde_\lambda(t) := \sum_{T \in \SYT(\lambda)} t^{\mathrm{maj}(T)}$ is the fake degree polynomial of $\lambda$ (equivalently, the graded multiplicity of the irreducible $S_n$-module $V_\lambda$ in the coinvariant algebra of $S_n$). By convention $\qe^{(0)} = 1$.

In \cite{Cli26-support} we proved the \emph{support theorem}: writing $\qe^{(n)} = \sum_\mu c_\mu^{(n)}\, p_\mu$ in the power-sum basis,
\[
  c_\mu^{(n)} \ne 0 \iff \mu = (e^{k}, \nu), \quad k := \lfloor n/e\rfloor, \quad \nu \vdash r_0 := n \bmod e,
\]
and moreover, with the notation $\alpha_r := 1 - \zetae^r$ and $b_r(n) := k + \one[r \le r_0]$,
\begin{equation}\label{eq:cnu-closed}
  c_{(e^k, \nu)}^{(n)} \;=\; \frac{1}{e^k\, z_\nu}\, \prod_{r=1}^{e-1} \alpha_r^{\,b_r(n) - m_r(\nu)},
\end{equation}
where $m_r(\nu)$ is the multiplicity of $r$ in $\nu$ and $z_\nu = \prod_i i^{m_i(\nu)}\, m_i(\nu)!$.

Since the parts of $\nu$ are $\le r_0 < e$, $p_{(e^k, \nu)} = p_e^{k}\, p_\nu$, so the support theorem factorises $\qe^{(n)}$ as
\[
  \qe^{(n)} \;=\; p_e^{k}\, F_e^{(n)}, \qquad F_e^{(n)} \;:=\; \sum_{\nu \vdash r_0} c_{(e^k, \nu)}^{(n)}\, p_\nu \;\in\; \Lambda_{r_0} \otimes \ZZ[\zetae].
\]
The main question left open by \cite{Cli26-support} was the structure of the ``residue factor'' $F_e^{(n)}$.

\begin{theorem}[Self-similarity]\label{thm:self-sim}
For every $n \ge 0$ and $e \ge 2$, with $k = \lfloor n/e\rfloor$ and $r_0 = n \bmod e$,
\[
  \qe^{(n)} \;=\; p_e^{\,k}\, \cdot\, \qe^{(r_0)}.
\]
Equivalently, $F_e^{(n)} = \qe^{(r_0)}$; in particular the residue factor depends on $n$ only through $r_0$, and its Schur expansion is
\[
  F_e^{(n)} \;=\; \sum_{\lambda \vdash r_0} \ftilde_\lambda(\zetae)\, s_\lambda.
\]
\end{theorem}

\section{Proof}

The proof is a direct calculation at the level of power-sum coefficients, using the closed form \eqref{eq:cnu-closed} and the classical cyclotomic identity below.

\begin{lemma}\label{lem:cyclo}
$\displaystyle \prod_{r=1}^{e-1} (1 - \zetae^r) \;=\; e.$
\end{lemma}

\begin{proof}
Setting $x = 1$ in the factorisation $x^{e} - 1 = (x - 1)\prod_{r=1}^{e-1}(x - \zetae^r)$ and taking the limit,
\[
  \prod_{r=1}^{e-1}(1 - \zetae^r) \;=\; \lim_{x \to 1}\, \frac{x^e - 1}{x - 1} \;=\; e. \qedhere
\]
\end{proof}

\begin{proof}[Proof of Theorem \ref{thm:self-sim}]
If $r_0 = 0$: then $b_r(n) = k$ for all $r \in [1, e-1]$, and $\nu$ ranges only over the empty partition, so $F_e^{(n)}$ is the constant symmetric function
\[
  F_e^{(n)} \;=\; \frac{1}{e^k}\, \prod_{r=1}^{e-1} \alpha_r^{\,k} \;=\; \frac{1}{e^k}\, e^k \;=\; 1
\]
by Lemma \ref{lem:cyclo}, and $\qe^{(0)} = 1$ by convention, so the theorem holds.

Now assume $r_0 \ge 1$. Apply the same closed form \eqref{eq:cnu-closed} at parameter $r_0$ (in place of $n$): then $\lfloor r_0/e\rfloor = 0$ (since $0 < r_0 < e$), and $b_r(r_0) = 0 + \one[r \le r_0] = \one[r \le r_0]$. So for every $\nu \vdash r_0$,
\begin{equation}\label{eq:cnu-r0}
  c_\nu^{(r_0)} \;=\; \frac{1}{z_\nu}\, \prod_{r=1}^{e-1} \alpha_r^{\,\one[r \le r_0] - m_r(\nu)}.
\end{equation}

Meanwhile from \eqref{eq:cnu-closed} at parameter $n$, with $b_r(n) = k + \one[r \le r_0]$,
\begin{align}
  c_{(e^k, \nu)}^{(n)}
  &= \frac{1}{e^k\, z_\nu} \prod_{r=1}^{e-1} \alpha_r^{\,k + \one[r \le r_0] - m_r(\nu)} \notag \\
  &= \frac{1}{e^k\, z_\nu}\, \Bigl(\prod_{r=1}^{e-1} \alpha_r\Bigr)^{k}\, \prod_{r=1}^{e-1} \alpha_r^{\,\one[r \le r_0] - m_r(\nu)} \notag \\
  &= \frac{1}{e^k\, z_\nu}\, e^{k}\, \prod_{r=1}^{e-1} \alpha_r^{\,\one[r \le r_0] - m_r(\nu)} \tag{by Lemma \ref{lem:cyclo}} \\
  &= c_\nu^{(r_0)}. \tag{by \eqref{eq:cnu-r0}} \label{eq:c-match}
\end{align}

Therefore
\begin{align*}
  \qe^{(n)}
  &= \sum_{\nu \vdash r_0} c_{(e^k, \nu)}^{(n)}\, p_{(e^k, \nu)} \tag{support theorem} \\
  &= \sum_{\nu \vdash r_0} c_{(e^k, \nu)}^{(n)}\, p_e^{k}\, p_\nu \\
  &= p_e^{k}\, \sum_{\nu \vdash r_0} c_\nu^{(r_0)}\, p_\nu \tag{by \eqref{eq:c-match}} \\
  &= p_e^{k}\, \qe^{(r_0)}. \tag{support theorem at $r_0$: all $\nu \vdash r_0$ have parts $< e$}
\end{align*}
The last step uses that when $r_0 < e$ the support theorem at parameter $r_0$ has $\lfloor r_0/e\rfloor = 0$, so \emph{all} partitions $\nu \vdash r_0$ contribute (their parts are automatically not divisible by $e$).

The Schur-basis statement follows: $F_e^{(n)} = \qe^{(r_0)} = \sum_{\lambda \vdash r_0} \ftilde_\lambda(\zetae) s_\lambda$ by definition of $\qe^{(r_0)}$.
\end{proof}

\section{Consequences}

\begin{corollary}[$n$-independence]\label{cor:n-ind}
The residue factor $F_e^{(n)} := \qe^{(n)}/p_e^{\lfloor n/e\rfloor}$ depends on $(n, e)$ only through $(e, n \bmod e)$.
\end{corollary}

\begin{proof}
Theorem \ref{thm:self-sim} identifies $F_e^{(n)}$ with $\qe^{(n \bmod e)}$, which manifestly depends only on $(e, n \bmod e)$.
\end{proof}

\begin{corollary}[Multiplicative recursion]\label{cor:recursion}
The sequence $\bigl(\qe^{(n)}\bigr)_{n \ge 0}$ satisfies
\[
  \qe^{(n)} \;=\; p_e\, \cdot\, \qe^{(n - e)} \qquad \text{for all } n \ge e.
\]
Equivalently, $\qe^{(qe + r)} = p_e^{q}\, \qe^{(r)}$ for $q \ge 0$ and $0 \le r < e$.
\end{corollary}

\begin{proof}
Direct from Theorem \ref{thm:self-sim}: $\qe^{(n)} = p_e^{\lfloor n/e\rfloor} \qe^{(n \bmod e)}$ and $\qe^{(n-e)} = p_e^{\lfloor n/e\rfloor - 1} \qe^{(n \bmod e)}$ have the same second factor.
\end{proof}

\begin{corollary}[Schur closed form for $F_e$]\label{cor:schur-closed}
For every $n$ with $r_0 = n \bmod e$,
\[
  F_e^{(n)} \;=\; \sum_{\lambda \vdash r_0} \ftilde_\lambda(\zetae)\, s_\lambda.
\]
In particular the Schur coefficient of $s_\lambda$ in $F_e^{(n)}$ equals $\ftilde_\lambda(\zetae)$: the fake degree of $\lambda$ evaluated at $\zetae$.
\end{corollary}

Corollary \ref{cor:schur-closed} explains the empirical pattern observed in the residue-factor probes of 2026-08-06:

\begin{itemize}
  \item \textbf{Boundary shapes have pure-root-of-unity coefficients.}
  For $\lambda = (r_0)$: $\ftilde_{(r_0)}(t) = 1$, so the coefficient of $s_{(r_0)}$ is $1$.

  For $\lambda = (1^{r_0})$: $\SYT(1^{r_0})$ has a unique tableau $T$ with $\mathrm{maj}(T) = \binom{r_0}{2}$, so $\ftilde_{(1^{r_0})}(t) = t^{r_0(r_0-1)/2}$, and the coefficient of $s_{(1^{r_0})}$ is $\zetae^{r_0(r_0-1)/2}$.

  \item \textbf{Interior shapes have cyclotomic-unit magnitudes.}
  For $\lambda = (2,1) \vdash 3$: $\SYT(2,1) = \bigl\{[1\,2;3],\, [1\,3;2]\bigr\}$ with major indices $\{2, 1\}$ respectively, so $\ftilde_{(2,1)}(t) = t + t^2 = t(1+t)$. At $\zetae$: $\zetae\,(1 + \zetae)$, with $|1+\zetae| = 2\cos(\pi/e)$. For $e = 5$ this gives the golden ratio $\varphi = 2\cos(\pi/5)$; for $e = 7$ it gives $2\cos(\pi/7)$; for $e = 4$ it gives $\sqrt 2$. Each of these values matches the empirical probe.
\end{itemize}

\begin{remark}[Semantic content of the theorem]
Theorem \ref{thm:self-sim} says that at the point $t = \zetae$ on the unit circle, the ``mass'' of the fake-degree Schur sum concentrates entirely on the coinvariant-algebra shape $(e^k, \nu)$ dictated by the arithmetic of $n$ modulo $e$; and, more surprisingly, the constant of proportionality is itself the same sum truncated to partitions of $r_0 = n \bmod e$. There is nothing new happening as $n$ increases past the ``fundamental domain'' $\{0, 1, \ldots, e-1\}$: the sequence $\qe^{(n)}$ is a $p_e$-linear extension of its first $e$ terms.
\end{remark}

\begin{remark}[The residue factor is really a $q_e$ in miniature]
Corollary \ref{cor:schur-closed} shows that $F_e^{(n)}$ is not merely a small polynomial in the roots of unity $\zetae^j$; it has the same intrinsic character as $q_e$ itself, only at parameter $r_0 < e$. The recursion in Corollary \ref{cor:recursion} then makes the entire family $\{\qe^{(n)}\}_{n \ge 0}$ into a $\ZZ[\zetae]$-lattice generated by the $e$ ``basic'' sums $\qe^{(0)}, \qe^{(1)}, \ldots, \qe^{(e-1)}$ under multiplication by $p_e$.
\end{remark}

\begin{remark}[Comparison with Kra\'skiewicz--Weyman]
Kra\'skiewicz--Weyman's theorem \cite{KW01} gives $\ftilde_\lambda(\zetae) = 0$ whenever $\lambda$ has non-empty $e$-core (in the generic Springer-regularity case). At first glance this suggests the Schur expansion of $q_e^{(r_0)}$ should be supported on partitions with empty $e$-core --- but for $\lambda \vdash r_0 < e$, the $e$-core condition is automatic (every partition of size $< e$ has $e$-core equal to itself, hence non-empty!) so KW does \emph{not} produce vanishing here. Instead all $\lambda \vdash r_0$ contribute with generically nonzero coefficient $\ftilde_\lambda(\zetae)$, which is a cyclotomic integer of degree $\le n(\lambda) \le \binom{r_0}{2} < \binom{e}{2}$. The KW vanishing is invisible below the ``horizon'' $r_0 < e$.
\end{remark}

\begin{remark}[Comparison with the Verschiebung route]
The 2026-08-05 dream cycle projected that Albion's Verschiebung theorem \cite{Alb26} would be the tool that closes the residue-factor question, via $\varphi_e s_\lambda = 0$ unless $e$-core$(\lambda) = \emptyset$. Theorem \ref{thm:self-sim} bypasses Verschiebung entirely: the collapse is a one-line consequence of the cyclotomic identity in Lemma \ref{lem:cyclo}, once the closed-form coefficient formula \eqref{eq:cnu-closed} is available (which itself comes from Chevalley--Molien, again bypassing Verschiebung). Both the support theorem and the self-similarity theorem are ``coinvariant algebra + roots of unity'' arithmetic, not Adams/Verschiebung representation theory.
\end{remark}

\section{Computational Verification}\label{sec:verif}

Independent numerical verification is at \verb|~/projects/probes/2026-08-07-self-similarity/verify.py|, testing 22 pairs $(n, e)$ at 50 decimal digits of precision:
\[
  (n, e) \in \{(2,3), (3,3), (4,3), (5,3), (6,3), (7,3), (8,3), (5,2), (6,2), (7,4),
\]
\[
  (10,4), (11,3), (7,5), (9,5), (9,4), (11,4), (11,5), (13,5), (12,6), (13,6), (15,4), (20,9)\}.
\]
For each pair we check two things, both by computing $c_\mu^{(n)} = z_\mu^{-1}\, \Psi_\mu(\zetae)$ directly from the coinvariant-algebra polynomial $\Psi_\mu(t) := \prod_{i=1}^n(1-t^i)/\prod_{j \in \mu}(1-t^j)$ (extracted as a polynomial in $\ZZ[t]$ via exact long division, then evaluated at the numerical $\zetae$):

\begin{enumerate}
  \item \textbf{Support:} $|c_\mu^{(n)}| < 10^{-30}$ for every $\mu \vdash n$ that is not of the form $(e^{\lfloor n/e\rfloor}, \nu)$ with $\nu \vdash r_0$.
  \item \textbf{Coefficient identity:} $|c_{(e^k, \nu)}^{(n)} - c_\nu^{(r_0)}| < 10^{-25}$ for every $\nu \vdash r_0$.
\end{enumerate}

All 22 pairs pass both tests. Sample:
\begin{center}
\begin{tabular}{rrrr r r}
\hline
$n$ & $e$ & $k$ & $r_0$ & off-supp max & $|c^{(n)}_{(e^k,\nu)} - c^{(r_0)}_\nu|$ max \\
\hline
$5$  & $3$ & $1$ & $2$ & $2.5 \times 10^{-51}$ & $2.4 \times 10^{-51}$ \\
$11$ & $3$ & $3$ & $2$ & $8.3 \times 10^{-51}$ & $1.6 \times 10^{-50}$ \\
$13$ & $5$ & $2$ & $3$ & $3.6 \times 10^{-51}$ & $4.7 \times 10^{-51}$ \\
$15$ & $4$ & $3$ & $3$ & $1.2 \times 10^{-50}$ & $2.3 \times 10^{-50}$ \\
$20$ & $9$ & $2$ & $2$ & $6.0 \times 10^{-50}$ & $6.3 \times 10^{-50}$ \\
\hline
\end{tabular}
\end{center}
The residuals are consistent with 50-digit round-off in $\zetae^i$ evaluations. All zero-support witnesses are numerical zeros; all coefficient matches agree beyond $10^{-25}$.

Verbose sample at $(n, e) = (13, 5)$, $k = 2$, $r_0 = 3$:
\begin{center}
\begin{tabular}{lll}
\hline
$\nu \vdash 3$ & $c_{(5,5,\nu)}^{(13)}$ & $c_\nu^{(3)}$ \\
\hline
$(3)$      & $+0.230328 - 0.708876\,i$ & $+0.230328 - 0.708876\,i$ \\
$(2,1)$    & $+0.904508 + 0.293893\,i$ & $+0.904508 + 0.293893\,i$ \\
$(1,1,1)$  & $-0.134836 + 0.414983\,i$ & $-0.134836 + 0.414983\,i$ \\
\hline
\end{tabular}
\end{center}

\section{Application: The Full $q_e$ Family Is a $p_e$-Lattice}

Corollary \ref{cor:recursion} gives $\qe^{(n)} = p_e^{\lfloor n/e\rfloor}\, \qe^{(n \bmod e)}$. So knowing the $e$ ``basic'' values
\[
  \qe^{(0)},\, \qe^{(1)},\, \ldots,\, \qe^{(e-1)}
\]
determines the entire sequence $\{\qe^{(n)}\}_{n \ge 0}$. Each basic value $\qe^{(r)}$ for $0 \le r < e$ is itself a small explicit symmetric function of degree $r$, with Schur coefficients $\ftilde_\lambda(\zetae) \in \ZZ[\zetae]$ ($\lambda \vdash r$).

For example, at $e = 3$:
\begin{align*}
  \qe^{(0)} &= 1, \\
  \qe^{(1)} &= s_{(1)}, \\
  \qe^{(2)} &= s_{(2)} + \zeta_3\, s_{(1,1)},
\end{align*}
and every other $\qe^{(n)}$ at $e = 3$ is $p_3^{\lfloor n/3\rfloor}$ times one of these three.

\section{Impact on the Composite-$d$ Program}

Combined with the support theorem \cite{Cli26-support} and Theorem B of \cite{Cli26-composite-d}, Theorem \ref{thm:self-sim} completes the character-level story of the residue factor $F_e$. In particular:

\begin{enumerate}
  \item \textbf{Explicit divisibility.} For any $(r, k, e)$ where Theorem B's factorisation-count obstruction fires, $\Phi_e(t) \mid m_\pi(t)$ for all $\pi \vdash r$, and the ``leftover'' coefficient factor is precisely the appropriately-scaled coefficient of $s_\lambda$ in $F_e$, which by Corollary \ref{cor:schur-closed} is $\ftilde_\lambda(\zetae)$ for $\lambda \vdash r_0$. There is nothing hidden.

  \item \textbf{No new tools needed.} The residue factor is not a mysterious cyclotomic-integer table awaiting classification: it \emph{is} a small copy of $q_e$ itself. Every question about $F_e$ reduces to the same question about $q_e^{(r_0)}$ with $r_0 < e$, where only $p(r_0)$ Schur summands appear.

  \item \textbf{The multiplicative structure is intrinsic.} Both the support theorem and the self-similarity theorem are consequences of the same underlying object: the graded character of the $S_n$-coinvariant algebra evaluated at $\zetae$. This is essentially the Reiner--Stanton--White CSP polynomial at $\zetae$, and the theorem crystallises the recursive structure that CSP polynomials have at roots of unity of order dividing the degree of the group.
\end{enumerate}

\begin{thebibliography}{9}
\bibitem{Cli26-support} Clio, \emph{Support of the fake-degree Schur sum at a root of unity}, \texttt{proofs/2026-08-06-support-conjecture-qe.tex}.
\bibitem{Cli26-composite-d} Clio, \emph{Composite-$d$ divisibility for wreath byproduct polynomials}, \texttt{proofs/2026-08-05-composite-d-k2-and-negative.tex}.
\bibitem{KW01} W. Kra\'skiewicz and J. Weyman, \emph{Algebra of coinvariants for a reflection group}, J. Algebra \textbf{240} (2001), 745--762.
\bibitem{Alb26} K. Albion, \emph{Character-factorisations and Verschiebung operators on symmetric functions}, arXiv:2501.18520 (to appear, Combinatorial Theory).
\bibitem{RSW04} V. Reiner, D. Stanton, D. White, \emph{The cyclic sieving phenomenon}, J. Combin. Theory Ser. A \textbf{108} (2004), 17--50.
\bibitem{Hum90} J. E. Humphreys, \emph{Reflection Groups and Coxeter Groups}, Cambridge Univ. Press, 1990.
\end{thebibliography}

\end{document}
